% Bu dosya iki biçimde kullanılabilir:
% 1) Tek başına derlenebilir.
% 2) book.tex içinden \subfile ile çağrılarak kitabın parçası olur.
\documentclass[00_kitap.tex]{subfiles}

\begin{document}
	\BolumKaynakcaBaslat

\subsubsection{Bir Foton Demetinin Kırınımı}

Genişliği \(a\) olan bir yarıktan geçen, dalga boyu \(\lambda\) olan
bir foton demetini ele alalım. Demeti önce elektromanyetik bir dalga
olarak düşünelim. Tek yarık kırınımından, yarıktan çıkan demetin
karakteristik açısal yayılmasının

\begin{equation}
	\theta \sim \frac{\lambda}{a}
\end{equation}

mertebesinde olduğu bilinir. Daha ayrıntılı kırınım hesabında ilk
minimum için

\begin{equation}
	\sin\theta_{\min}=\frac{\lambda}{a}
\end{equation}

elde edilir.

Şimdi aynı olayı fotonların parçacık özelliği açısından ele alalım.
Yarıktan önce fotonların yaklaşık olarak \(x\)-doğrultusunda ilerlediğini
düşünelim. Bir fotonun momentum büyüklüğü

\begin{equation}
	p=\frac{h}{\lambda}
\end{equation}

olduğundan, başlangıçta

\begin{equation}
	p_x\simeq p=\frac{h}{\lambda},
	\qquad
	p_y\simeq 0
\end{equation}

olarak düşünülebilir.

Foton yarıktan geçtiğinde, \(y\)-doğrultusundaki konumu yarık
genişliğiyle sınırlanır. Dolayısıyla fotonun bu doğrultudaki konumunun
karakteristik belirsizliği

\begin{equation}
	\Delta y \sim a
\end{equation}

mertebesindedir.

Konum ile momentum arasındaki belirsizlik bağıntısı

\begin{equation}
	\Delta y\,\Delta p_y
	\geq
	\frac{\hbar}{2}
\end{equation}

olduğundan, \(y\)-doğrultusundaki konumun sınırlandırılması,
\(p_y\) momentum bileşeninin artık tek bir değerle belirlenememesi
anlamına gelir. Karakteristik olarak

\begin{equation}
	\Delta p_y
	\sim
	\frac{\hbar}{a}
\end{equation}

mertebesinde bir enine momentum yayılması beklenir.

Böylece yarıktan sonra fotonların momentumlarının büyüklüğü yaklaşık
olarak aynı kalırken, momentum vektörlerinin yönleri farklılaşabilir.
Küçük açılar için

\begin{equation}
	\theta
	\sim
	\frac{p_y}{p_x}
\end{equation}

olduğundan, açısal yayılmanın karakteristik büyüklüğü

\begin{equation}
	\theta
	\sim
	\frac{\Delta p_y}{p_x}
\end{equation}

şeklinde tahmin edilebilir.

Burada

\begin{equation}
	p_x\simeq \frac{h}{\lambda}
\end{equation}

olduğuna göre,

\begin{equation}
	\theta
	\sim
	\frac{\hbar/a}{h/\lambda}
	=
	\frac{1}{2\pi}\frac{\lambda}{a}.
\end{equation}

Bu hesapta ortaya çıkan sayısal katsayıya özel bir anlam
yüklenmemelidir. Çünkü burada \(\Delta y\), \(\Delta p_y\) ve
\(\theta\) yalnızca karakteristik genişlikleri temsil etmektedir;
örneğin \(\Delta p_y\), kırınım deseninin ilk minimumuna karşılık gelen
momentum değeri değildir.

Belirsizlik ilkesi bu nedenle kırınım açısının tam değerini değil,
kırınımın karakteristik ölçeğini verir:

\begin{equation}
	\boxed{
		\theta \sim \frac{\lambda}{a}
	}.
\end{equation}

İlk minimumun tam konumu ise elektromanyetik dalga kuramından veya
eşdeğer olarak fotonun dalga fonksiyonunun Fourier dönüşümünden

\begin{equation}
	\boxed{
		\sin\theta_{\min}=\frac{\lambda}{a}
	}
\end{equation}

olarak elde edilir.

Böylece dalga ve parçacık betimlemeleri birbirleriyle tutarlı bir
sonuca ulaşır: fotonun yarıktan geçerken \(y\)-doğrultusundaki
konumunun sınırlandırılması, \(y\)-doğrultusundaki momentum
bileşenlerinin yayılmasına; bu da yarıktan sonra fotonların farklı
açılarda gözlenebilmesine karşılık gelir.










\BolumKaynakcaBitir
\end{document}
